% Chapter 5: Proof Sketch of Mihalcea's Positivity Theorem
% Draft version — to be integrated into chapter 5
% Based on Mihalcea (2007) \cite{Mihalcea}

\section{Proof of Mihalcea's Positivity Theorem}
\label{sec:MihalceaProof}

In this section we give a sketch of the proof of Mihalcea's main positivity theorem,
referring to \cite[Section 6, pages 8--11]{Mihalcea} for full details.

\subsection{Main Theorem}

\begin{Theorem}[Mihalcea Positivity Theorem {\cite[Theorem 1]{Mihalcea}}]
\label{th:MihalceaPositivity}
Let $u, v, w \in W^{P}$ and let $d$ be a multidegree.
Then the equivariant quantum Littlewood--Richardson coefficient
$c_{u,v}^{w,d}$ belongs to $\Lambda = \mathbb{Z}[x_{1},\ldots,x_{r}]$
and is a polynomial with nonnegative integer coefficients.
\end{Theorem}

The proof proceeds by reducing the computation of $c_{u,v}^{w,d}$ to
Graham's positivity theorem on the product space $X \times X$.
This reduction is the content of \S 6 of \cite{Mihalcea}.

\subsection{Setup}

Let
\begin{itemize}
\item $X = G/P$ be a homogeneous space,
\item $B = T \cdot U$ the Borel subgroup, $B^{-}=w_{0}Bw_{0}$ the opposite Borel,
\item $B' = T\cdot(U^{-}\times U^{-})$ the semidirect product,
\item $\overline{\mathcal{M}}=\overline{\mathcal{M}}_{0,3}(X,d)$ the stable map moduli space,
\item $ev_{i}:\overline{\mathcal{M}}\to X$ the evaluation maps ($i=1,2,3$),
\item $F=(ev_{1},ev_{2}):\overline{\mathcal{M}}\to X\times X$,
\item $\Lambda=H_{T}^{*}(pt)=\mathbb{Z}[x_{1},\ldots,x_{r}]$.
\end{itemize}

The torus $T$ acts on $X\times X$ via the diagonal action.
The semidirect product $B'$ acts on $X\times X$ by
\[
t\cdot(u_{1},u_{2})(\bar g_{1},\bar g_{2})
=(tu_{1}\bar g_{1},\,tu_{2}\bar g_{2}).
\]

The equivariant quantum Littlewood--Richardson coefficient is defined by
\begin{equation}
c_{u,v}^{w,d}
=\pi_{*}^{T}\bigl((ev_{1}^{T})^{*}\sigma(u)^{T}
\cdot (ev_{2}^{T})^{*}\sigma(v)^{T}
\cdot (ev_{3}^{T})^{*}\widetilde{\sigma}(w)^{T}\bigr),
\label{eq:EQCLRDef}
\end{equation}
where $\pi:\overline{\mathcal{M}}\to\operatorname{Spec}\Lambda$ is the projection
and $\widetilde{\sigma}(w)^{T}$ denotes the equivariant quantum Schubert class
(Poincar\'e dual of $\sigma(w)^{T}$).

\subsection{Three auxiliary results}

The proof rests on three geometric ingredients from \cite{Mihalcea}.

\begin{Lemma}[{\cite[Lemma 6.4]{Mihalcea}}; $ev_{3}^{-1}(Y(w))$ components]
\label lem:ev3inverse}
Let $w\in W^{P}$. Then $ev_{3}^{-1}(Y(w))$ is $T$--stable, reduced,
and is a disjoint union of irreducible components $V_{i}$, each of codimension
$\ell(w)$ in $\overline{\mathcal{M}}$.
\end{Lemma}

\begin{Lemma}[{\cite[Lemma 6.3]{Mihalcea}}; Kleiman transversality]
\label lem:KleimanTransversality}
Let $Z$ be an irreducible $G$--variety and
$F:Z\to X\times X$ a $G$--equivariant morphism
($G$ acts diagonally on $X\times X$).
Then for any $u,v\in W^{P}$,
$F^{-1}(X(u)\times Y(v))$ is reduced and locally irreducible,
of codimension $c(u)+\ell(v)$ in $Z$.
\end{Lemma}

\begin{Corollary}[{\cite[Corollary 6.2]{Mihalcea}}; Graham positivity on $X\times X$]
\label cor:GrahamOnXX}
Let $V\subset X\times X$ be a $T$--stable subvariety.
In $H_{T}^{*}(X\times X)$ we have
\[
[V]_{T}=\sum_{i}f_{i}\,[Y(u_{i})\times Y(v_{i})]_{T},
\qquad f_{i}\in\Lambda,
\]
where each $f_{i}$ is a polynomial with nonnegative integer coefficients.
\end{Corollary}

\subsection{Key pushforward formula}

From Lemma \ref{cor:GrahamOnXX} we obtain the decomposition
\[
(ev_{3}^{T})^{*}[Y(w)]_{T}=\sum_{i}a_{i}\,[V_{i}]_{T},
\qquad a_{i}\in\mathbb{N},
\tag{6.5}
\]
where $a_{i}$ is the degree of the restricted map $F|_{V_{i}}:V_{i}\to F(V_{i})$.

Applying the pushforward formula to $F$ and using the projection formula
$f_{*}^{T}((f^{T})^{*}(a)\cdot b)=a\cdot f_{*}^{T}(b)$ (see \cite[(1)]{Mihalcea}),
equation \eqref{eq:EQCLRDef} becomes
\[
c_{u,v}^{w,d}
=\sum_{i}a_{i}\,
(\pi_{X\times X})_{*}^{T}
\bigl([X(u)\times X(v)]_{T}\cdot[F(V_{i})]_{T}\bigr).
\tag{6.6}
\tag{$*$}
\end{equation}
This is the \textbf{key pushforward formula}.

\subsection{Proof of Theorem \ref{th:MihalceaPositivity}}

\begin{Proof}[Sketch]
Starting from \eqref{eq:EQCLRDef}, we apply the three steps of \S 6.2 in
\cite{Mihalcea}:

\textbf{Step 1.}  Use Lemma \ref lem:ev3inverse} to decompose the factor
$(ev_{3}^{T})^{*}\widetilde{\sigma}(w)^{T}$ as in (6.5).

\textbf{Step 2.}  Apply the projection formula and the push--forward formula
to obtain the key pushforward formula ($*$).

\textbf{Step 3.}  For each irreducible component $V_{i}$, consider its image
$F(V_{i})\subset X\times X$.  By Corollary \ref cor:GrahamOnXX} we may write
\[
[F(V_{i})]_{T}=\sum_{j}f_{j}^{i}\,[Y(u_{j})\times Y(v_{j})]_{T},
\qquad f_{j}^{i}\in\Lambda\ \text{nonnegative}.
\]

Substituting into ($*$) and using the projection formula once more, we obtain
\[
c_{u,v}^{w,d}
=\sum_{i}\sum_{j}
a_{i}\,f_{j}^{i}\,
(\pi_{X\times X})_{*}^{T}
\bigl([X(u)\times X(v)]_{T}\cdot[Y(u_{j})\times Y(v_{j})]_{T}\bigr).
\]

Finally, by the equivariant Poincar\'e pairing
(see \cite[Lemma 2.3]{Mihalcea})
\[
(\pi_{X\times X})_{*}^{T}
\bigl([X(u)\times X(v)]_{T}\cdot[Y(u_{j})\times Y(v_{j})]_{T}\bigr)
=\delta_{u,u_{j}}\cdot\delta_{v,v_{j}},
\]
which is a nonnegative integer.

Since all coefficients $a_{i}$ are nonnegative integers and all $f_{j}^{i}$
are polynomials with nonnegative integer coefficients, the entire sum
$c_{u,v}^{w,d}$ is a polynomial in $\Lambda$ with nonnegative integer coefficients.
This is precisely the statement of Theorem \ref{th:MihalceaPositivity}.
\qedhere
\end{Proof}

\subsection{Remark on $T'$--equivariant positivity}

Mihalcea also proves a version for a larger torus $T'=(\mathbb{C}^{\times})^{m}$.
Since $T=T'/\mathbb{C}^{\times}$, the ring $H_{T'}^{*}(pt)=\mathbb{Z}[T_{1},\ldots,T_{m}]$
contains $H_{T}^{*}(pt)$ via the relations $T_{j}-T_{j+1}=x_{j}$.

\begin{Corollary}[{\cite[Corollary 7.1]{Mihalcea}}; $T'$--equivariant positivity]
\label cor:TdashPositivity}
In $QH_{T'}^{*}(X)$ the equivariant quantum Littlewood--Richardson coefficients
are polynomials in the differences $T_{1}-T_{2},\ldots,T_{m-1}-T_{m}$
with nonnegative integer coefficients.
\end{Corollary}

This corollary bridges the present result with the triple positivity framework
of \cite{GX2025}, where the positivity is expressed in the variables $t_{i}-y_{j}$.

\endinput
